% ==============================================================================
% conclusion.tex
% ==============================================================================

\section{Conclusion}
\label{sec:conclusion}

This survey has presented a systematic review of learning-based methods for few-view X-ray-to-CT reconstruction, covering the period from 2018 to 2025. We organised the literature into four major architectural families — GAN-based, transformer-based, neural field, and diffusion-based methods — and analysed their respective strengths, limitations, and trade-offs.

Our review highlights several key findings:

\begin{itemize}
  \item The field has evolved rapidly from early convolutional encoder--decoder networks to sophisticated generative models capable of producing high-fidelity 3D volumes from as few as a single X-ray.
  \item Diffusion-based methods represent the current state of the art in perceptual quality, though at significant computational cost.
  \item The domain gap between synthetic training data and real clinical X-rays remains the principal barrier to clinical translation.
  \item Standardised benchmarks and evaluation protocols are urgently needed to enable fair comparison and track genuine progress.
\end{itemize}

As deep generative models continue to advance and computational resources become more accessible, we anticipate that few-view X-ray-to-CT reconstruction will move closer to clinical viability, potentially transforming imaging workflows in radiation therapy planning, orthopaedic surgery, and emergency medicine.

% TODO: Finalise conclusion with additional reflections from the full review.
